\section{引言：为什么推理能力是 AI 的关键}

Denny Zhou 是 Google DeepMind 的研究科学家，长期专注于大语言模型的推理能力研究。他是 Chain-of-Thought Prompting 和 Self-Consistency 等里程碑工作的核心贡献者。本次讲座从一个根本性问题出发——\textbf{什么才算真正的智能？}——引出 LLM 推理的核心脉络。

\begin{knowledgebox}{从机器学习到推理的范式转变}
传统机器学习（few-shot learning、active learning、meta-learning 等）追求数据高效性（data efficiency），但在实践中这些方法并没有真正实现"从少量样本学习"。Denny Zhou 认为，人类能够从少量样本中学习，根本原因不是统计规律，而是\textbf{推理能力（reasoning）}。这一洞察驱动了他从机器学习转向 LLM 推理的研究方向。
\end{knowledgebox}

\subsection{Last Letter Concatenation：一个简单但深刻的例子}

Denny Zhou 用 Last Letter Concatenation 任务来揭示传统方法与推理方法的巨大差异：

\begin{itemize}
    \item \textbf{任务定义}：给定一个人名（如 ``Elon Musk''），输出名和姓最后一个字母的拼接（``n'' + ``k'' = ``nk''）。
    \item \textbf{传统 ML 方法}：用 Transformer 训练，需要数千标注样本，准确率约 85--90\%。
    \item \textbf{Few-shot Prompting}：直接给 LLM 几个示例，LLM 仍然容易出错（如 ``Barack Obama'' 输出 ``ck'' 而非 ``ka''）。
    \item \textbf{Chain-of-Thought Prompting}：在示例中加入推理步骤（``The last letter of Elon is n, the last letter of Musk is k, so nk''），仅用一个示例即可达到 100\% 准确率。
\end{itemize}

\begin{importantbox}{核心洞察}
对于人类来说简单的任务，如果一个方法需要海量标注数据才能学会，那它难以称为真正的「智能」。推理能力使 LLM 能够从极少样本中泛化，这是传统机器学习方法无法比拟的。
\end{importantbox}

\subsection{本章小结}

推理能力是弥合「数据驱动学习」与「真正智能」之间鸿沟的关键。LLM 通过 Chain-of-Thought 式的中间步骤生成，实现了质的飞跃——从需要大量训练数据降低到仅需少量甚至零样本。

